\section{System Design}
\label{sec:design}

\system makes two selection decisions before placing replicas: which
Agents' histories deserve protection (O1), and which KV regions within
those histories should be retained (O2). It then places the selected
regions in Peer HBM or DRAM, adapting the Peer budget to foreground
demand (O3). Both tiers hold partial backups. The design accounts for
the same Agent reusing its history across successive invocations, as
well as for the workflow delay caused by losing that history.

\subsection{Overview and Invocation State}
\label{sec:teacher-overview}

Figure~\ref{fig:teacher-overview} shows the control and data paths.
The workflow runtime supplies invocation dependencies, progress, and
accepted token histories. The \emph{Workflow Tracker} estimates remaining
work and workflow slack. The \emph{Recovery-aware Router} dispatches
invocations using load and valid prefix locality, and identifies
compatible takeover workers. The \emph{KV Master} selects Agent histories
and KV regions, then assigns their replicas to memory. A per-worker
\emph{Adapter} publishes the available Peer budget; rank-local
\emph{Connectors} execute Master's transfers and report completion.

\begin{figure}[t]
 \centering
 \includegraphics[width=\columnwidth]{figures/teacher-system-overview.pdf}
 \caption{Selection followed by two-tier placement. DAG nodes are Agent
 invocations: A(1) and A(2) invoke the same Agent. Each DP worker has two TP
 ranks; selected Peer regions map to corresponding ranks on a compatible
 worker. DRAM holds selected regions assigned to it. Blue paths carry
 KV; dashed paths carry control information. Backups and token histories
 used for recovery must survive the protected source failure.}
 \label{fig:teacher-overview}
 \Description{Application, Controller, and Engine bands. A(1) forks into
 B(1) and C(1), which join at A(2). Tracker and Router feed KV Master. Master
 selects Agents and regions before Connectors place replicas. Adapter
 reports Peer capacity. Two DP workers contain TP0 and TP1 execution
 and Peer pools, with rank-matched copies and partial DRAM backups.}
\end{figure}

\paragraph{One node is one invocation.}
A DAG node $v=(w,a,r)$ identifies the $r$th invocation of Agent $a$ in
workflow $w$. Thus, $a^{(1)}$ and $a^{(2)}$ are distinct nodes even when they
execute consecutively and use the same model. The later invocation may
extend the earlier input and accepted outputs. Its reusable KV is the
materialized state for the exact common token prefix, subject to matching
model revision, positional metadata, numerical format, and KV layout.
Agent identity alone does not establish compatibility; an edited or
truncated prompt can shorten the reusable prefix. A dependency between
different Agents likewise does not imply that they share KV.

Figure~\ref{fig:teacher-invocation-reuse} illustrates this distinction.
KV lifetime follows the history and its consumers, rather than the
completion of one invocation. Master keeps a per-history manifest and
consumer references, so a sealed region produced by $a^{(1)}$ can serve
$a^{(2)}$ and $a^{(3)}$ without allocating a new backup for each call. Completing
$a^{(1)}$ updates its remaining consumers and value; it does not automatically
retire its replicas. If future calls are not yet visible, a configured
idle-history window permits retention subject to capacity pressure.
Unknown future calls receive no invented DAG nodes or guaranteed value.
Consumer references inform retention value; only active physical readers
and transfers prevent immediate reclamation.

\begin{figure}[t]
 \centering
 \includegraphics[width=\columnwidth]{figures/teacher-invocation-reuse.pdf}
 \caption{Successive invocations of one Agent reuse growing history.
 In this append-only example, A(2) consumes the compatible prefix from
 A(1), and A(3) consumes the history extended by A(2). Each box in the top
 row is a separate DAG node. Lower bars show logical token histories,
 not complete backup allocations; only selected KV regions are retained.}
 \label{fig:teacher-invocation-reuse}
 \Description{Three successive invocation nodes A(1), A(2), and A(3) for
 Agent A. A(2)'s input contains history H1 plus new input U2; A(3)'s input
 contains H1, the extension E2 produced by A(2), and new input U3.
 H1 and H1 plus E2 are the compatible reusable prefixes.}
\end{figure}

\paragraph{Selected state and its location.}
Let $\mathcal K_a$ be the sealed, materialized KV regions of a compatible
Agent history, and $S_a\subseteq\mathcal K_a$ its selected backup set.
Each region identifies a layer and a historical token block together
with its required TP shards. The stable placement is
\begin{equation}
 S_a=H_a\mathbin{\dot\cup}D_a,\qquad
 H_a\subseteq\{b\in S_a:\operatorname{layer}(b)\le L\},
 \label{eq:teacher-placement}
\end{equation}
where $H_a$ and $D_a$ reside in Peer HBM and DRAM, respectively. The
disjoint union denotes one managed backup per selected region; source
execution KV is separate. Neither tier, nor their union, is required
to cover the complete history. Temporary migration copies and all TP
allocations consume physical capacity until released. Compatible
physical objects referenced by multiple consumers are charged once.

Backup destinations lie outside the protected source failure domain.
For host-failure protection, this also excludes DRAM on the source host.
The token log and replica directory must survive that failure. These
requirements provide the inputs for restore or reconstruction, extending
checkpoint-and-replay recovery~\cite{strati2024dejavu,cao2026lumen}; they
do not turn a partial manifest into a complete checkpoint.

\subsection{O1: Selecting Agent Histories}
\label{sec:teacher-agent-protection}

O1 selects histories according to their effect on the remaining workflow.
Tracker maintains a residual DAG of individual invocations. Completed
nodes contribute no remaining execution, but their retained histories
can still benefit later nodes. Running nodes contribute only unfinished
work. Consecutive calls of the same Agent remain separate dependent
nodes, linked to their common history manifest.

\paragraph{Estimating urgency.}
A GBDT estimates remaining invocation duration from currently available
features: model and Agent role, input length, decoded-token progress,
elapsed execution, and completed calls' lengths and timings. Training
and calibration separate whole workflows. Forward and backward
critical-path passes~\cite{kelley1959cpm} produce earliest and latest
finishes, $EF_v$ and $LF_v$, and guarded slack
\begin{equation}
 s_v=\max\{0,LF_v-EF_v-g_v\},
 \label{eq:teacher-slack}
\end{equation}
where $g_v$ is a calibrated upper quantile of slack overestimation.
Successive invocations share path slack; it cannot be spent independently
at every node. Stale or invalid estimates grant no positive slack.
Slack orders urgent consumers, while the joint schedule below evaluates
the effect of protecting their histories.

\paragraph{Reuse and recovery value.}
For a selected set and placement $P$, let $F_w^0(P)$ be the predicted
finish of workflow $w$ during normal execution, and $F_w^f(P)$ its finish
after a common failure of source worker $f$. Normal execution includes
the Router's actual prefix locations, expected execution-KV retention,
and replica-copy interference. A backup earns no normal-reuse benefit
when the same prefix remains available in execution HBM. If the Router
uses a backup for a later invocation, its transfer and any missing-state
reconstruction are prerequisites of that invocation.

The failure schedule includes all interrupted invocations and later
visible consumers of history lost with the worker. Restoration,
dependency reconstruction, and accepted-token replay precede unfinished
execution. Takeover admission and transfers contend on the resources
they actually share. A fixed ready-time rule, using guarded slack and
request ID to break ties, applies to every candidate. Master compares
\begin{equation}
 J(P)=(1-\lambda)\sum_w F_w^0(P)
       +\lambda\sum_f\pi_f\sum_w F_w^f(P),
 \label{eq:teacher-objective}
\end{equation}
where $0\le\lambda\le1$ is a configured recovery weight and $\pi_f$ are
fixed normalized source-failure weights, equal by default. These are
policy weights, not measured fault probabilities. The comparison fixes
the workload snapshot, Router targets, weights, and queue rules.

A single physical history can accelerate several later invocations,
but each invocation appears only once in each schedule. The objective
therefore avoids adding a full-history benefit for every reference.
Normal prefix reuse is evaluated separately from the fault-time slack
of the currently running call. At a join, the maximum predecessor
finish determines progress, preserving O1's non-additive effect.

\paragraph{Generating Agent candidates.}
Master first identifies Agent histories with active or visible future
consumers. Candidates include single-history changes and small bundles
covering co-bottleneck paths into an exposed join. Protecting only one
of two equally delayed branches may have zero effect, while protecting
both can advance the join. Each Agent candidate passes to O2 to
select useful regions; selection does not reserve a full checkpoint.
Master evaluates the resulting candidate, placement, and victims
together. Bounded search can miss other combinations and does not claim
global optimality. Without a valid DAG profile, it uses profiled
invocation time saved per byte with deterministic ties, retaining the
same correctness and capacity rules.

\subsection{O2: Selecting Regions, Then Placing Replicas}
\label{sec:teacher-region-selection}

Within each candidate history, O2 selects the KV whose retention can
reduce exposed waiting or reconstruction. The output is a logical set
$S_a$ independent of its eventual memory addresses. Only after forming
this set does Master assign its members to Peer HBM or DRAM. Physical
feasibility can require revising the selection, but there is no initial
step that copies the complete history into DRAM.

\paragraph{Selecting KV regions.}
For region $b$ in layer $\ell$, recent completed native attention
observations supply
\begin{equation}
 \widehat p_{a,\ell,b}=
 \frac{\text{eligible selection units that selected }b}
      {\text{selection units for which }b\text{ was visible}}.
 \label{eq:teacher-demand}
\end{equation}
The unit follows the model's native selection rule; execution tiling
does not redefine it. Head duplicates count once per unit, and causal
visibility is determined before pruning. Small denominators use a
calibrated layer prior. Comparable visible blocks in dense attention
have access factor one. Frequency orders region candidates; it neither
guarantees future demand nor establishes a time saving.

Master considers individual regions and progressively larger hot sets,
including dependency bundles needed to make a proposed restoration
useful. Recent consecutive demand windows preserve correlated access
and native reveal order. Their size and the profiling work budget are
fixed parameters. Invalid sparse statistics use a conservative full
causal demand envelope for profiling, without forcing a full backup.
Neither selection nor placement changes native attention semantics.

\paragraph{Cost of a missing region.}
The profiler distinguishes a valid local execution hit, a Peer hit, a
DRAM fetch, and a region with no surviving copy. For the last case,
the engine reconstructs the required dependency closure from durable
tokens and reusable state. Rebuilding a block's upper-layer KV may
require lower-layer activations and causally preceding KV; the cost is
not an independent one-block Prefill. Native selector and positional
state are also dependencies. Shared reconstruction tasks are executed
and charged once within a recovery schedule.

If the engine cannot safely combine partial state with reconstruction,
it prefills the complete required prefix and gives those unusable
partial replicas no avoided-compute credit. A partial backup therefore
does not promise to avoid full reconstruction. Predictions can affect
latency, but every native demand must be satisfied before attention
continues. This is the correctness path for both later invocations
and failure recovery.

\paragraph{Measuring exposed delay.}
The profile schedules copies, bounded prefetches, dependency computation,
and unfinished invocation work on a common timeline. A demand-driven
fetch cannot begin before the selector reveals it. If $q_\ell$ is the
time a layer is ready to consume history, $a_\ell(P)$ the time its last
required region becomes available, and $c_\ell$ its remaining compute,
then
\begin{equation}
 \begin{aligned}
 e_\ell(P)&=[a_\ell(P)-q_\ell]^+,\\
 f_\ell(P)&=q_\ell+e_\ell(P)+c_\ell.
 \end{aligned}
 \label{eq:teacher-overlap}
\end{equation}
Availability may require a transfer or reconstruction. Query and
selector preparation, resource queues, and dependencies determine
readiness; later layers depend on earlier finishes. Only computation
that remains to execute can hide a transfer. The resulting schedule
feeds Equation~\ref{eq:teacher-objective}; historical frequency is not
used as a substitute for exposed delay.

\paragraph{Placing the selected set.}
The first $L$ model layers are eligible for Peer HBM, where $L$ is a
configured depth for a compatible model and layout. Master first offers
selected regions from these layers to Peer. When they exceed its quota,
it orders them by the profiled incremental workflow benefit of Peer
over DRAM per byte, retaining incumbents on ties. Each admission must
fit every required TP rank. Selected early-layer regions that do not
fit go to DRAM, together with selected regions beyond layer $L$.
Thus, $L$ limits Peer placement; it does not select all KV in those
layers or require complete later layers in DRAM.

DRAM also has a finite budget. If the selected set cannot fit across
the two tiers, Master removes or replaces eligible regions in increasing
recomputed workflow loss per released byte and reevaluates the whole
plan. Regions that remain unbacked use the reconstruction path.
Optional replacements require $J(P)-J(P')>\epsilon$ and feasible
transient memory and transfer reservations. Ranking is recomputed
after tentative changes, including joint Agent bundles and victims.
Per-round work and copy-byte limits bound overhead and churn. Selection
and placement are distinct decisions within this bounded search;
placement costs are included before committing either.

\begin{figure}[t]
 \centering
 \includegraphics[width=\columnwidth]{figures/teacher-selective-recovery.pdf}
 \caption{Partial selection before placement, illustrated with $L=2$.
 Five of twelve logical regions are selected. Peer holds L1/B0 and
 L1/B2; L2/B1 falls back to DRAM because the illustrative Peer quota is
 full. Selected L3/B0 and L3/B2 also go to DRAM. Unselected cells have
 no managed backup. A demand uses a surviving copy or reconstructs its
 missing dependencies. Each cell denotes all required TP shards;
 this is a constructed state, not a timing result.}
 \label{fig:teacher-selective-recovery}
 \Description{A selected-region mask followed by an exclusive HBM/DRAM
 placement map. Two selected first-layer regions reside in Peer, an
 overflowing second-layer region resides in DRAM, and two third-layer
 regions reside in DRAM. Seven regions are unselected. The two tiers
 together store five regions, not a complete twelve-region checkpoint.}
\end{figure}

\subsection{O3: Adapting the Peer HBM Budget}
\label{sec:teacher-capacity}

Master chooses contents and locations; Adapter determines the Peer
capacity foreground execution can support. For takeover rank $k$,
$B^{\rm pool}_k$ excludes weights and fixed allocations,
$B^{\rm serve}_k$ covers the serving high watermark and execution
reservations, and $B^{\rm reserve}_k$ reserves growth and workspace.
The ceiling is
\begin{equation}
 C_k=\max\{0,\min(B^{\rm cap}_k,
 B^{\rm pool}_k-B^{\rm serve}_k-B^{\rm reserve}_k)\}.
 \label{eq:teacher-ceiling}
\end{equation}
Peer use is not part of $B^{\rm serve}_k$ and is not subtracted twice.
Committed, reserved, copying, and retiring allocations all count toward
physical occupancy. Distinct TP ranks have distinct ceilings.

Adapter uses additive increase and multiplicative
decrease~\cite{chiu1989aimd}. Allocation stalls, reserve violations, or
memory-related admission waits multiply quota $B_k$ by $0<\gamma<1$.
After $m$ stable windows below a lower wait threshold, Adapter adds
$\Delta$ bytes if Master reports beneficial early-layer candidates
currently in DRAM or deferred for Peer capacity. The signal is computed
before quota filtering, allowing growth from zero. Updates are
page-aligned and clipped to $C_k$; missing telemetry disables growth.
Adapter changes total quota, while the configured depth $L$ remains
unchanged.

When quota shrinks, Master stops new Peer admissions and moves the
least valuable Peer placements to DRAM if their regions remain selected.
It reserves DRAM space, copies and validates the destination, redirects
the manifest, then releases the Peer allocation after readers drain.
There is no assumed DRAM backing to discard against. If DRAM cannot
accommodate the migration, Master reselects the finite backup set,
explicitly removing regions and their protection credit as needed.
Their token histories still permit reconstruction. Mandatory pressure
relief need not improve $J$, but cannot bypass active physical leases.
Foreground admission uses actual occupancy; a reduced quota alone
does not free memory. Pending copies may be cancelled, and allocations
wait when immediate physical space is unavailable.

\subsection{Online Replication and Recovery}
\label{sec:teacher-execution}

Planning runs on invocation arrival and completion, history or DAG
updates, memory pressure, and a fixed block-aligned interval $Z$ of
accepted tokens. Prefill and decode offer immutable materialized regions
as they become available. Final invocation completion offers its sealed
state and updates future-consumer references. The interval controls
review frequency; it is not a complete-checkpoint or replay-tail bound.

Algorithm~\ref{alg:teacher-protection} applies O1 and O2 before assigning
memory locations. Every candidate comparison includes current replicas,
victims, migrations, and normal-service copy interference. Version checks
cover the DAG, history, manifest, and quotas. A stale or infeasible
plan is requeued without advertising its proposed replicas. Candidate
gain is normalized by changed destination bytes for ordering; each
tier's physical capacity is enforced separately.

\par\medskip\noindent\begin{minipage}{\columnwidth}
\raggedright
\refstepcounter{designalgorithm}\label{alg:teacher-protection}
\textbf{Algorithm \thedesignalgorithm: Select, place, and reconcile}
\par\smallskip
\begin{tabularx}{\columnwidth}{@{}r@{\quad}X@{}}
\toprule
1 & Reconcile copies; snapshot invocations and history references \\
2 & $P\gets$ current placement; apply mandatory capacity relief \\
3 & $A\gets\textsc{AgentCandidates}(P)$, including joint bundles \\
4 & \textbf{while} planning work budget remains \textbf{do} \\
5 & \quad Form candidate sets $S_c\gets\textsc{SelectRegions}(A,P)$ \\
6 & \quad For each $c$, form $P_c\gets\textsc{Place}(S_c,L,\text{budgets})$ \\
7 & \quad Fit victims and re-place if a budget is exceeded \\
8 & \quad Reprofile each whole plan, including missing-state work \\
9 & \quad $P^*\gets$ best feasible gain/byte candidate \\
10 & \quad \textbf{if} no gain exceeds $\epsilon$ \textbf{then break} \\
11 & \quad $P\gets P^*$; refresh affected candidates \\
12 & \textbf{if not} \textsc{ReserveIfCurrent}$(P)$ \textbf{then} \\
13 & \quad Requeue event; \textbf{return} \\
14 & Copy selected regions directly to their assigned tier \\
15 & Publish completed objects; retire old owners after drain \\
\bottomrule
\end{tabularx}
\end{minipage}\par\medskip

\paragraph{Publishing and retaining partial state.}
A manifest records each committed region's token interval, layer,
required shards, compatibility metadata, and current owner. Connector
publishes a region only after all its required objects are validated;
pending copies cannot satisfy a lookup. For migration, old and new
allocations are both charged until the destination commits and old
readers drain. Failed copies leave the previous committed owner valid
when it survives. Retention follows live history references and current
selection value, including successors of a completed invocation.

The highest saved token position is not a complete checkpoint frontier.
Only an explicitly verified all-layer, all-shard prefix may be treated
as complete. Sparse manifests preserve holes, and the planner includes
their reconstruction cost. With no active or visible future consumer,
an idle history expires at its configured retention limit; workflow
completion releases its remaining logical references. Physical reuse
by another compatible consumer or an existing reader delays release.

\paragraph{Consuming retained history.}
On a normal subsequent invocation, Router matches its exact token prefix
against execution state and surviving manifests. It uses compatible
local KV first, then available replicas, and reconstructs missing state.
An earlier call's output can enter this prefix only where its KV has
actually been materialized and validated.

On confirmed failure, Router fences the old execution epoch and selects
a compatible surviving target. Connector leases the surviving regions;
the engine combines valid hits, DRAM reads, and dependency reconstruction
without omitting native attention demands. Peer residency counts as a
local hit only on its actual target; another target must transfer that
state or reconstruct it. The engine replays recorded accepted tokens
where necessary and resumes generation without emitting them twice.
Full-prefix Prefill remains the fallback when partial restoration is
unsupported or no useful state survives. These mechanisms specify a
design; integrated measurements are still needed to establish recovery
gains, cross-invocation reuse, and normal-service overhead.
